\chapter*{Abstract}
\addcontentsline{toc}{chapter}{Abstract}

Most Applicant Tracking Systems (ATS) today rely on rigid keyword matching, often rejecting highly qualified candidates simply because they use different words to describe their experience. To solve this problem, I developed an automated, human-centric resume screening pipeline using Natural Language Processing (NLP). Instead of merely counting overlapping words with traditional sparse arrays (like TF-IDF), my research uses advanced semantic embeddings (Sentence-BERT) to understand the actual meaning behind a candidate's skills. Beyond text matching, I also sought to capture a human recruiter's intuition regarding a candidate's tenure. To achieve this, I mathematically modeled the ``Experience Gap'' ($\Delta E$) by comparing a piecewise quadratic penalty against a novel Sigmoid-Bayesian framework. This Bayesian approach effectively mimics human judgment, bypassing the need for massive historical hiring datasets that are rarely available.

To ensure my models could genuinely generalize to new, unseen job postings without memorizing specific vocabularies, I evaluated the entire pipeline using a strict \texttt{GroupShuffleSplit} architecture. I benchmarked traditional regression and binary classification models against specialized Learning-to-Rank (LTR) algorithms. The results empirically demonstrate that resume screening should not be treated as a simple pass/fail classification, but rather as a listwise sorting problem. Ultimately, my best-performing model --- an \texttt{LGBMRanker} powered by LambdaMART gradients and dense semantic embeddings --- achieved an outstanding top-10 ranking accuracy (NDCG@10 of $0.9398$). These findings confirm that embracing listwise ranking algorithms over rigid keyword filters can fundamentally transform and improve how automated systems identify top talent.